\chapter{相关技术与理论基础}
\vspace{7pt}

本章对自动化用例建模系统所涉及的关键技术进行系统梳理，为后续章节的系统设计与实现提供理论基础。首先阐述以Transformer为核心的大语言模型原理，重点分析DeepSeek推理模型在结构化输出方面的技术特性；随后介绍检索增强生成技术中文本向量化与语义匹配的基本原理；接着讨论多智能体系统的角色分工理论以及LangGraph基于状态图的工作流编排机制；最后介绍PlantUML建模语言与Vue 3前端框架的集成技术。

\section{大语言模型原理}

\subsection{Transformer架构与注意力机制}

Vaswani等人\cite{vaswani2017attention}于2017年首次提出的Transformer架构，是当前大语言模型的核心基础。该架构完全摒弃了此前序列建模中广泛使用的循环神经网络和卷积神经网络，转而采用自注意力机制作为序列信息提取的主要手段。自注意力机制的核心思想在于，对于输入序列中的每一个词元，模型通过计算其与序列中所有其他词元的注意力权重，动态聚合全局上下文信息。

自注意力机制的计算过程可以形式化描述如下。给定输入序列的嵌入表示矩阵X，模型首先通过三个可学习的权重矩阵W\_Q、W\_K和W\_V将其分别映射为查询矩阵Q、键矩阵K和值矩阵V。注意力权重通过Q与K的点积并经过缩放和Softmax归一化得到，最终的输出为注意力权重与V的加权求和。这一计算过程可表示为公式(2.1)。

% 公式2.1 自注意力机制计算
\begin{equation}
    \text{Attention}(Q, K, V) = \text{Softmax}\left(\frac{QK^T}{\sqrt{d_k}}\right)V
    \label{eq:attention}
\end{equation}

其中，$d_k$表示键向量的维度，除以$\sqrt{d_k}$起到缩放作用，防止点积值过大导致Softmax函数进入梯度饱和区域。通过自注意力机制，模型能够在单层内直接建立序列中任意两个位置之间的依赖关系，而无需像循环神经网络那样通过逐步传递的方式，这显著提升了长距离依赖的建模能力。

Transformer架构采用编码器-解码器的结构形式。编码器由多个相同的层堆叠而成，每层包含一个多头自注意力子层和一个前馈全连接子层，每个子层后均接有残差连接和层归一化。多头注意力机制通过并行运行多组独立的注意力计算，使模型能够从不同表示子空间中提取信息，进一步增强了模型的表达能力。前馈全连接子层则由两个线性变换和一个非线性激活函数组成，对每个位置的表示进行独立处理。解码器的结构与编码器类似，但增加了一个交叉注意力子层，用于在生成过程中关注编码器的输出表示。

基于Transformer架构的大语言模型通过在超大规模文本语料上进行预训练，学习到了丰富的语言知识和语义表示能力。预训练通常采用自监督学习方式，常见的训练目标包括掩码语言模型和自回归语言模型两种范式。通过在海量数据上进行预训练，模型不仅习得了词汇、语法等表层语言知识，还获得了常识推理、逻辑判断等深层次语义理解能力，这为其在下游任务中的迁移应用奠定了能力基础。

\subsection{DeepSeek推理模型及其在结构化JSON输出中的优势}

DeepSeek是由深度求索公司研发的大语言模型系列，本研究采用DeepSeek系列的推理增强模型作为核心智能引擎。DeepSeek模型在Transformer架构的基础上进行了多项改进优化。在注意力机制方面，模型采用了多头潜在注意力机制，通过对键值对进行低秩压缩，在保持长上下文建模能力的同时显著降低了推理过程中的内存占用。在混合专家架构方面，模型将前馈网络层替换为多个专家子网络，每个词元根据路由策略仅激活其中部分专家进行计算，以较小的计算代价实现了模型总容量的扩展。Nair等人\cite{nair2025finetuned}在微调大语言模型与传统NLP管道的对比研究中，也证实了DeepSeek系列模型在结构化文本生成任务上的出色表现。

在自动化用例建模系统中，大语言模型的一个核心任务是从自然语言需求文本中提取结构化信息，如参与者列表、用例名称、实体类及其属性与方法等。模型输出的信息需要以机器可解析的结构化格式呈现，以便后续工作流节点进行处理。JSON因其简洁清晰的键值结构成为本系统中模型输出的主要格式。然而，在实际工程实践中，大语言模型生成的JSON文本经常出现格式不合法的现象，常见的错误类型包括缺少闭合括号、键名未加引号、数组元素间遗漏逗号以及在JSON结构前后附加非JSON的说明性文本。

DeepSeek推理模型在结构化JSON输出方面展现出若干优势。模型在预训练和微调阶段接触了大量结构化数据，对JSON格式的语法约束具有较强的内化理解，能够在大多数情况下输出合法的JSON字符串。更为关键的是，DeepSeek推理模型采用了思维链式的推理增强策略\cite{deepseek2025r1}：在处理复杂任务时，模型会在内部进行逐步的逻辑推理，梳理清楚实体之间的关系和各字段的填充依据，然后再输出最终的结构化结果。这种先推理后输出的机制有效降低了由于推理跳跃导致的字段遗漏和格式错乱问题。

本系统在后端实现中进一步引入了JSON输出的容错处理机制。当模型输出的字符串无法被标准JSON解析器直接解析时，系统会尝试进行格式修复，包括截取第一个合法JSON对象、补齐缺失的闭合括号以及清除附加的非JSON文本。这种双重的质量保障策略使系统在实际运行中保持了较高的结构化输出成功率。

\section{检索增强生成技术}

\subsection{文本向量化与相似度匹配原理}

检索增强生成（Retrieval-Augmented Generation，RAG）技术是一种将信息检索与文本生成相结合的技术范式，旨在增强大语言模型在知识密集型任务中的表现\cite{brown2025systematic}。大语言模型虽然通过预训练习得了广泛的知识，但其知识截止于训练数据的收集时间，且在面对需要精确引用特定文档内容的任务时容易产生幻觉。RAG技术通过在生成之前先从外部知识库中检索与当前问题相关的文档片段，并将检索结果作为上下文信息嵌入到提示词中，从而引导模型基于可靠信源进行生成。

文本向量化是实现语义检索的基础技术环节，其目标是将不定长的自然语言文本映射到固定维度的向量空间中的点，使得语义相似的文本对应的向量在空间中彼此接近。现代文本向量化通常采用预训练的句子嵌入模型。本研究选用的嵌入模型为paraphrase-multilingual-MiniLM-L12-v2，这是一个在多语言场景下经过优化的轻量级模型，能够较好地处理中英文混合文本的语义表示。

文本向量之间的相似度通常采用余弦相似度进行度量。给定两个文本的向量表示a和b，其余弦相似度定义为两向量内积除以它们的模长乘积，如公式(2.2)所示。

% 公式2.2 余弦相似度
\begin{equation}
    \text{cosine\_similarity}(a, b) = \frac{a \cdot b}{\|a\| \cdot \|b\|}
    \label{eq:cosine}
\end{equation}

余弦相似度的取值范围在-1到1之间，数值越接近1表示两个向量方向越一致，即对应文本的语义越相近。在检索实践中，系统将查询文本向量化后，计算其与知识库中所有文档块向量的余弦相似度，按相似度由高到低排序，选取排名靠前的文档块作为检索结果。与关键词匹配检索相比，这种基于语义的检索方式能够发现表述方式不同但语义相关的内容，具有更强的泛化能力。

\subsection{向量数据库ChromaDB在需求上下文检索中的作用}

ChromaDB是一个专为向量检索场景设计的嵌入式数据库，本系统选择其作为RAG模块的向量存储与检索引擎。与其他向量数据库相比，ChromaDB的显著特点是部署轻量、接口简洁，能够以嵌入式模式直接集成到Python应用中，无需额外的服务进程或复杂的集群配置。Nguyen等人\cite{nguyen2026novel}在自动化UML图生成统一框架中同样采用了ChromaDB作为RAG的向量存储引擎，验证了其在需求上下文检索场景中的适用性。

在本系统的RAG架构中，ChromaDB承担了向量索引存储和相似度检索两类核心职责。系统在文档处理阶段将需求文本按段落或固定长度窗口分块，调用嵌入模型将每个文本块转化为向量，然后将向量连同原始文本和元数据一起存入ChromaDB的集合中。当用户启动对某个模块的用例提取任务时，系统使用该模块的需求描述文本作为查询，经同样的嵌入模型向量化后，向ChromaDB发起相似度检索请求。ChromaDB在内部维护了高效的向量索引结构，能够在大规模向量集合中快速定位与查询向量最相似的若干个向量，并将对应的原始文本块作为检索结果返回。

检索返回的上下文文本块随后被拼接到发送给大语言模型的提示词中。通过这种方式，模型在生成时不仅依据用户直接输入的模块需求描述，还能参考需求文档中其他相关段落的完整语义信息。在长文档场景下，这一机制尤为重要。以大篇幅的需求规格说明书为例，某个模块的核心功能描述可能集中在文档的某一章节，但其前置条件、异常流程和接口约束等信息可能散布在文档的其他章节中。若仅将当前模块的局部文本输入模型，模型可能因缺少全局信息而遗漏关键实体或关系。RAG通过语义检索将这些分散的相关信息汇集到提示词中，使模型获得更完整的信息视图，从而生成更准确的用例模型。

\section{多智能体系统与LangGraph}

\subsection{智能体角色分工理论}

多智能体系统是指由多个自主的智能体组成的系统，这些智能体通过协作共同完成单个智能体难以独立处理的复杂任务。其中，任务分解、角色分工和通信协作是多智能体系统的三大核心机制。在多智能体架构中，角色分工是提升系统整体效能的关键设计原则，其优势体现在多个方面。

角色分工使每个智能体的策空间缩小，降低了单个智能体的认知负荷和决策复杂度。在用例建模任务中，如果将参与者提取和关系分析等全部交给同一个智能体完成，该智能体需要同时关注自然语言中的多个维度的信息，容易出现注意力分散和信息遗漏。针对智能Agent的感知、决策与执行等特性进行系统化的需求建模与模式划分，已被证明能够有效指导复杂智能系统的工程化设计\cite{xiao2026problemframes, hong2024metagpt, li2026aiagent}。

专业化分工还使得各智能体的输出格式能够针对其子任务进行定制化设计。参与者提取智能体的输出只需是一个参与者名称列表，关系分析智能体的输出则是一个关系三元组列表。这种针对性的输出设计不仅使每个智能体的生成目标更加明确，也便于下游节点进行解析和校验。

此外，多智能体架构天然支持递进式求解。复杂任务往往具有内在的步骤依赖性，如关系分析必须在实体提取完成的基础上进行。多智能体架构通过将此类依赖关系映射为智能体之间的调用顺序，使任务的执行流程与问题本身的逻辑结构保持一致，有助于提升求解质量。

\subsection{基于状态图的工作流循环与中断机制}

LangGraph是一个面向大语言模型应用的工作流编排框架，其核心概念是将多步骤的智能体协作流程建模为有向状态图。在LangGraph的抽象中，工作流中的每个处理步骤被定义为一个节点，节点之间通过有向边连接，表示数据和执行权的流转方向。所有节点共享一个全局的状态对象，节点在执行时从状态对象中读取输入数据，并将处理结果写回状态对象。这种基于状态图的设计使复杂工作流的定义变得直观且易于维护。在工业界的复杂数据流处理与人工智能建模平台中，类似的有向图调度机制已被证明能显著降低开发门槛并提升任务执行效率\cite{xu2025cloudnative}。

状态图模型具有较强的表达能力，能够描述丰富的工作流控制模式。顺序执行是最基本的模式，即节点A的输出直接传递至节点B。条件分支允许根据状态对象中的某些字段值决定下一步应激活哪个节点，这在需要根据中间结果动态调整流程时十分实用。循环模式允许在满足特定条件时将执行权重新导向到某上游节点，从而支持迭代优化的处理逻辑。这些控制模式的组合为构建灵活的智能体协作流程提供了坚实的基础。

在自动化用例建模系统中，LangGraph的状态图承载了整个建模工作流。系统的状态对象UMLState包含了输入的需求文本、目标图表类型、已提取的实体列表、已分析的关系数据以及当前的工作状态标识等字段。工作流中的每个节点对应一个专门的建模智能体或处理函数。以用例图的生成为例，提取实体节点首先从需求文本中识别参与者和用例名称并写入状态，随后关系分析节点读取实体列表并分析关系类型。当需要人工确认时，流程在特定的节点之前暂停，等待用户通过前端界面确认或修改中间结果后再继续执行。

中断机制是多智能体协作系统中实现人机协同的关键技术手段。LangGraph提供了在指定节点执行前暂停工作流的能力，称为断点。当工作流到达设置了断点的节点时，状态图的执行被挂起，当前的状态对象被持久化保存。此时用户可以通过应用界面查看当前的中间结果，对提取出的实体列表进行修改或确认。用户提交确认后，工作流从断点处恢复执行，继续后续节点的处理。这种机制使用例建模过程不再是一个完全自动化的黑箱，而是变成了人机交互协作的增量式建模过程，既发挥了自动化处理的效率优势，又保留了人类专家对关键决策的控制权。

\section{建模语言与可视化技术}

\subsection{PlantUML语法规范}

PlantUML是一种基于文本描述的UML图表生成工具，用户通过编写简洁的文本语法描述模型结构，由PlantUML引擎自动渲染为高质量图形。与传统图形界面建模工具相比，其文本驱动的特点在多方面体现出优势。Cámara等人\cite{camara2023assessment}在评估ChatGPT建模性能时将PlantUML作为目标输出格式，指出文本形式的模型描述天然适配版本控制系统，支持差异比较、合并与历史追溯，这是二进制图形文件难以实现的功能。戴莉萍和王文乐\cite{dailiping2021}在UML建模实验中也指出，编程式建模在可维护性、协作效率和自动化程度方面均优于传统拖拽式建模。

从表达效率看，PlantUML以紧凑的文本语法替代了繁琐的图形操作。参与者、用例、类及其属性方法、消息序列等建模元素均可通过关键字和简洁的结构化语句声明，少量代码即能描述内容丰富的模型图，且渲染引擎自动完成布局计算，免去了手动调整元素位置的工作。从可扩展性看，PlantUML在核心语法之外还提供了构造型标注、皮肤参数自定义等扩展机制，用户可根据需要定制图形的视觉风格和语义标注。此外，PlantUML的文本语法结构规整、关键词明确，使其生成任务易于与大语言模型对接——大语言模型只需掌握有限的语法规则即可输出合法代码，这一特性是本系统选择PlantUML作为输出格式的重要考量。在建模过程中，规范的建模语法表达是实现自动化建模的基础保障。

\subsection{Vue3响应式框架与Monaco Editor集成技术}

前端可视化技术在本系统中承担着用户交互与图表展示的重要职责。系统前端采用Vue 3框架构建。Vue 3引入了组合式应用程序接口（Composition API），通过setup函数和响应式引用机制提供了更灵活的组件逻辑组织方式。组合式API使开发者能够将相关联的状态和逻辑聚合在一起，形成可复用的组合函数，显著提升了代码的可维护性。在响应式原理方面，Vue 3基于JavaScript的代理对象实现了细粒度的数据响应式追踪，只有实际被访问的数据属性才会被追踪，从而提升了大型应用中的性能表现。这些特性对于本系统中多个数据表格与图表组件之间需要保持状态同步的复杂交互场景尤为重要。

Monaco Editor是微软开源的浏览器端代码编辑器，也是Visual Studio Code的核心编辑组件。本系统在前端中集成Monaco Editor作为PlantUML代码的编辑和展示工具。Monaco Editor提供了语法高亮、代码折叠、自动补全和错误提示等专业级代码编辑功能，能够为用户提供与桌面集成开发环境接近的编辑体验。

本系统对Monaco Editor的集成不仅限于基本的代码编辑，还实现了编辑区与数据表单之间的双向同步功能。当用户在表格中修改参与者名称或用例名称时，系统会重新调用Jinja2模板渲染PlantUML代码并更新编辑器内容。反之，当用户直接在编辑器中修改PlantUML代码时，系统通过解析器从代码中逆向提取出实体和关系数据，并同步更新表格显示。这种双向同步机制使用户可以选择自己偏好的交互方式进行建模操作，在可视化的表格操作和文本式的代码操作之间自由切换，提升了系统的灵活性和操作效率。

\section{本章小结}

本章以Transformer及注意力机制为起点，阐明了大语言模型处理自然语言的原理，并分析了DeepSeek在结构化输出方面的优势。针对需求文档信息分散问题，引入RAG与ChromaDB向量数据库，为模型提供上下文支持。为进一步提升复杂建模的可靠性与可控性，多智能体角色分工与LangGraph状态图工作流结合，实现了任务递进求解与人工中断确认的协同机制。最后，PlantUML的文本驱动表达与Vue 3集成Monaco Editor的前端方案，构成了可编辑、可追溯的建模交互环境。上述技术共同构成了系统设计与实现的理论基础，下一章将在此基础上有序展开需求分析。